problem:  
Let \( M^{2n} \) be a compact, simply connected, even-dimensional manifold admitting a Riemannian metric of almost positive sectional curvature—that is, \( K_p(\sigma) > 0 \) on an open dense subset of \( G_2(TM) \), the Grassmannian of tangent 2-planes.  Assume \( M \) also admits an isometric, effective torus action \( T^k \curvearrowright M \) of maximal possible rank permitted by its dimension (for instance, \( k = n - 1 \) or \( n - 2 \)).  

**Question:**  
Is it possible for such a torus-equivariant almost-positively curved manifold \( M^{2n} \) to have an *indefinite intersection form* on \( H^n(M;\mathbb{R}) \)?  Equivalently, does the presence of a large symmetry group together with almost positive curvature force the middle-dimensional cohomology to be definite?  

Why is it a "good" problem:  
This formulation probes the boundary between *curvature rigidity* and *topological freedom* in a precise and testable way.  Known examples of almost-positive curvature (such as Eschenburg spaces and Bazaĭkin biquotients) overwhelmingly occur in highly symmetric settings and in dimensions where the middle cohomology is trivial—suggesting a deep link between symmetry, topology, and lower curvature bounds.  

By introducing a large torus symmetry, one brings the problem into interaction with major structural tools from transformation group theory (e.g., the connectedness lemma, equivariant formality, and the work on positively curved manifolds with torus actions).  Determining whether almost positive curvature combined with such symmetry enforces definiteness of the intersection form would illuminate how far the classic topological restrictions of positive curvature extend once curvature positivity is relaxed on a set of measure zero.  

The problem is *simple to state* yet *deeply structural*.  It sits at the interface of Riemannian geometry, equivariant topology, and rational homotopy theory.  A resolution—either the discovery of an explicit counterexample or a rigidity mechanism ensuring definiteness—would significantly advance our understanding of how geometry and symmetry together determine global manifold topology.  Its narrow formulation around torus actions provides concrete leverage, making this an excellent and genuinely “good” research problem.